\begin{figure}[tb]
\centering
\begin{tikzpicture}[scale=.95]
 \begin{scope}[xshift=-3.3cm]
  \draw[->] (-2.3,0)--(2.4,0) node[right]{$i_d$};
  \draw[->] (0,-1.55)--(0,1.65) node[above]{$i_{\mathrm{exc}}$};
  \draw[loss curve] (0,0) ellipse (1.85 and .82);
  \node at (0,-1.95){amperes: unequal loss price};
 \end{scope}
 \draw[-{Latex[length=3mm]},very thick] (-.8,0)--(.8,0)
 node[midway,above]{$\matr R^{1/2}$};
 \begin{scope}[xshift=3.3cm]
  \draw[->] (-1.65,0)--(1.75,0) node[right]{$y_d$};
  \draw[->] (0,-1.55)--(0,1.65) node[above]{$y_e$};
  \draw[loss curve] (0,0) circle (1.15);
  \node at (0,-1.95){distance: equal loss price};
 \end{scope}
\end{tikzpicture}
\caption{Loss whitening changes units, not physics. The ellipsoid becomes a
sphere because one coordinate unit is chosen to cost one unit of copper loss
in every direction.}
\label{fig:whitening}
\end{figure}
